% Image Analysis Report Template
% Extends the core mycelium report template with imaging-specific sections.

\documentclass[11pt,a4paper]{article}

\usepackage[utf8]{inputenc}
\usepackage[T1]{fontenc}
\usepackage{amsmath,amssymb}
\usepackage{graphicx}
\usepackage{booktabs}
\usepackage{hyperref}
\usepackage[margin=1in]{geometry}
\usepackage{natbib}
\usepackage{xcolor}
\usepackage{caption}
\usepackage{subcaption}
\usepackage{siunitx}            % For consistent unit formatting

\title{Image Analysis: [Title]}
\author{[Author(s)]}
\date{\today}

\begin{document}

\maketitle

\begin{abstract}
[Summarize: biological question, imaging modality, number of images/conditions analyzed,
segmentation approach, key quantitative findings, main conclusion.]
\end{abstract}

\section{Introduction}

[Biological context. What was imaged and why? What quantitative measurements
are needed to answer the question?]

\section{Methods}

\subsection{Image Acquisition}

\begin{itemize}
    \item \textbf{Microscope}: [model]
    \item \textbf{Objective}: [magnification, NA]
    \item \textbf{Pixel size}: \SI{[X]}{\micro\meter/pixel}
    \item \textbf{Modality}: [widefield/confocal/light-sheet]
    \item \textbf{Channels}: [list with excitation/emission]
    \item \textbf{Exposure}: [time per channel]
\end{itemize}

\subsection{Image Preprocessing}

[Describe background subtraction, illumination correction, registration.
Document parameters for each step.]

\subsection{Segmentation}

[Describe segmentation method, software version, all parameters.
Reference validation results.]

% \begin{figure}[htbp]
%     \centering
%     \begin{subfigure}{0.32\textwidth}
%         \includegraphics[width=\textwidth]{../outputs/figures/raw_example.png}
%         \caption{Raw image.}
%     \end{subfigure}
%     \begin{subfigure}{0.32\textwidth}
%         \includegraphics[width=\textwidth]{../outputs/figures/segmentation_overlay.png}
%         \caption{Segmentation overlay.}
%     \end{subfigure}
%     \begin{subfigure}{0.32\textwidth}
%         \includegraphics[width=\textwidth]{../outputs/figures/segmentation_labels.png}
%         \caption{Label image.}
%     \end{subfigure}
%     \caption{Representative segmentation result.}
%     \label{fig:segmentation}
% \end{figure}

\subsection{Quantification}

[Describe features measured, normalization applied, physical units used.]

\subsection{Statistical Analysis}

[Describe statistical methods for comparing conditions. Reference core
statistical conventions. Document sample sizes at each level (images,
objects, biological replicates).]

\section{Results}

\subsection{Quality Control}

[Summarize image QC: number of images acquired, excluded, and why.
Show representative QC figures.]

\subsection{Segmentation Validation}

[Report segmentation performance metrics against ground truth if available.
Show representative good and challenging cases.]

\subsection{Quantitative Analysis}

[Report measurements and comparisons. Show violin/box plots for per-object
measurements. Report effect sizes and confidence intervals.]

% \begin{figure}[htbp]
%     \centering
%     \includegraphics[width=0.8\textwidth]{../outputs/figures/measurement_comparison.pdf}
%     \caption{[Description of quantitative comparison between conditions.]}
%     \label{fig:measurements}
% \end{figure}

\section{Discussion}

[Interpret results. Discuss segmentation limitations. Consider biological
implications. Propose follow-up experiments or analyses.]

\bibliographystyle{plainnat}
\bibliography{references}

\end{document}
